    \documentclass[11pt,a4paper]{article}

    % ---------- Simple, self-contained formatting ----------
    \usepackage[margin=2cm]{geometry}
    \usepackage{graphicx}
    \usepackage{tabularx}
    \usepackage{booktabs}
    \usepackage{array}
    \usepackage{float}
    \usepackage{xcolor}
    \usepackage{amsmath}
    \usepackage{amssymb}
    \usepackage{enumitem}
    \usepackage{microtype}
    \usepackage{tikz}   % [Draft 1] schematic figures
    \usetikzlibrary{arrows.meta,decorations.pathreplacing}

    \usepackage[numbers,sort&compress]{natbib}

    \usepackage[
        colorlinks=true,
        linkcolor=black,
        citecolor=blue,
        urlcolor=blue
    ]{hyperref}

    % ---------- [Draft 1] Change mark-up ----------
    % Everything added or altered in Draft 1 is shown in green.
    % For the final submission change \showchangestrue to \showchangesfalse:
    % the text stays and only the green colouring is removed.
    \newif\ifshowchanges
    \showchangestrue
    \definecolor{addgreen}{RGB}{0,135,60}
    \ifshowchanges
        \newcommand{\added}[1]{\textcolor{addgreen}{#1}}
        \newenvironment{addedblock}{\color{addgreen}}{}
        \newcommand{\addedfloat}{\color{addgreen}}  % floats do not inherit colour
    \else
        \newcommand{\added}[1]{#1}
        \newenvironment{addedblock}{}{}
        \newcommand{\addedfloat}{}
    \fi
    % [TBC] = data the group must supply or confirm before submission
    \newcommand{\tbc}[1]{\added{\textbf{[TBC\if\relax\detokenize{#1}\relax\else: #1\fi]}}}
    % References added in Draft 1 are also shown in green in the reference list
    \makeatletter
    \newcommand*{\addedrefs}[1]{\@for\@tempa:=#1\do{%
        \expandafter\gdef\csname addedref@\@tempa\endcsname{}}}
    \AtBeginDocument{\ifshowchanges
        \let\draft@lbibitem\@lbibitem
        \def\@lbibitem[#1]#2{\normalcolor\draft@lbibitem[#1]{#2}%
            \@ifundefined{addedref@#2}{}{\color{addgreen}}\ignorespaces}%
    \fi}
    \makeatother
    \addedrefs{campbell2004,astrom1997,hexcel2013,williams1996,twigg2004,explorecomposites_release,toray_t300,astm_d3171,krenchel1964,judd1978,centea2015,swolfs2014}
    % Ragged-right table column for the new tables
    \newcolumntype{L}{>{\raggedright\arraybackslash}X}



    \begin{document}

    % ======================================================
    % COVER PAGE
    % ======================================================
    \begin{titlepage}
        \centering
        \vspace*{\fill}
        \includegraphics[width=0.6\textwidth]{usyd_logo.png}\\[1em]
        {\scshape\Large School of Aerospace, Mechanical and Mechatronic Engineering}\\[3em]
        {\scshape\Large MECH2460/8660 Manufacturing Engineering}\\[1em]
        \rule{\linewidth}{0.5mm}\\[0.4cm]
        {\Huge\bfseries Polymer Composite Manufacturing Lab Report}\\[0.1cm]
        \rule{\linewidth}{0.5mm}\\[3em]
        {\large Canvas Group: 72}\\[0.5em]
        {\large Laboratory date: 08/10/2026}\\[3em]

        \renewcommand{\arraystretch}{1.5}

\begin{tabularx}{\linewidth}{|l|l|X|}
    \hline
    Name & SID & Contribution \\ \hline
    Abdullah Al Nuhian & 550232069 & Abstract, Introduction and Methodology \\ \hline
     & & \\ \hline
     & & \\ \hline
     & & \\ \hline
\end{tabularx}

        \vspace*{\fill}
    \end{titlepage}

    \pagenumbering{arabic}


    % =============================================
    % Main Document
    % =============================================

    \section{Abstract}

\begin{addedblock}
This report describes the manufacture and characterisation of a
fibre-reinforced polymer (FRP) aerofoil made by manual prepreg lay-up,
vacuum bagging and oven curing. Each skin was laid up in an aluminium
billet mould as a symmetric four-ply hybrid laminate, with woven
carbon/epoxy on both faces and two woven E-glass/epoxy plies at the
centre ($[\mathrm{C}/\mathrm{G}]_\mathrm{s}$), and the skins were
bonded to a polyurethane foam core. The theoretical overview covers
mould making (female and male moulds, process-specific requirements
and material selection, mould manufacturing methods, gelcoats and
release agents) and the role of each vacuum-bagging consumable. From
the measured skin thickness of \tbc{}\,mm and the nominal fabric areal
masses and fibre densities, the fibre volume fraction was estimated as
\tbc{}\,\%, against about 49\,\% expected for a void-free laminate; the
mould-side surfaces \tbc{one-line summary of surface quality}. The
discussion compares common FRP manufacturing processes and examines
how fibre type, volume fraction, orientation, stacking sequence, voids
and cure conditions govern laminate properties, showing that placing
the carbon plies outermost roughly doubles the flexural rigidity of
the skin. Small UAV wing skins are proposed as a suitable application
for hand lay-up.
\end{addedblock}



\section{Introduction}

\subsection{Fibre-Reinforced Polymer Composites}

Fibre-reinforced polymer (FRP) composites are multiphase engineering
materials comprising reinforcing fibres embedded within a continuous
polymer matrix. The fibres provide most of the load-bearing strength
and stiffness, while the matrix maintains fibre alignment, transfers
inter-fibre stresses and protects the reinforcement from environmental
damage \citep{mallick2007}. The resulting mechanical properties are
generally anisotropic, meaning they depend on loading direction.
Through appropriate selection of fibre type, orientation, volume
fraction and stacking sequence, FRP laminates can be tailored to
specific structural requirements.

A major advantage of FRPs is their high specific strength and
stiffness, enabling substantial weight savings in applications where
structural efficiency is critical. Unlike metals, which generally
exhibit ductile behaviour, or ceramics, which offer high hardness and
thermal resistance but are typically brittle, FRPs provide
directionally optimised mechanical performance and often excellent
corrosion resistance. However, limitations include relatively low
through-thickness strength, susceptibility to delamination and
temperature restrictions imposed by the polymer matrix
\citep{mallick2007}. Table~\ref{tab:material_comparison} summarises
the principal engineering differences between these material classes.

\begin{table}[htbp]
    \centering
    \caption{Qualitative comparison of FRP composites with
    conventional engineering materials \citep{mallick2007}.}
    \label{tab:material_comparison}
    \small
    \renewcommand{\arraystretch}{1.15}
    \begin{tabularx}{\linewidth}{@{}p{2.5cm}XXX@{}}
        \toprule
        \textbf{Property} &
        \textbf{FRP Composites} &
        \textbf{Metals} &
        \textbf{Ceramics} \\
        \midrule
        Specific strength
        and stiffness &
        Often high; direction-dependent &
        Moderate to high &
        Material-dependent \\

        Mechanical
        behaviour &
        Anisotropic; may delaminate &
        Generally ductile &
        Hard but brittle \\

        Corrosion
        resistance &
        Often good &
        Alloy-dependent &
        Generally good \\

        Temperature
        resistance &
        Limited by matrix &
        Alloy-dependent &
        Often excellent \\

        Manufacturing &
        Lay-up, consolidation and curing &
        Casting, forming and machining &
        Forming, sintering and finishing \\
        \bottomrule
    \end{tabularx}
\end{table}

These characteristics have established FRPs in aerospace, automotive,
marine and renewable-energy applications, where lightweight
construction and tailored structural performance are particularly
valuable \citep{wu2026}.

\subsection{Fibre Architectures and Matrix Materials}

Common reinforcing fibres include carbon, glass and aramid. Carbon
fibres offer high specific stiffness and strength, making them
suitable for lightweight structural components. Glass fibres
provide a cost-effective combination of strength, durability and
electrical insulation, while aramid fibres exhibit high toughness
and impact resistance \citep{mallick2007}.

Fibre architecture also strongly influences composite performance.
Unidirectional (UD) reinforcement concentrates strength and
stiffness along the fibre direction, whereas woven fabrics provide
reinforcement along multiple yarn directions. Plain-weave fabrics
have frequent interlacements, improving fabric stability, while
twill weaves generally offer better drapability over curved
surfaces due to reduced interlacement frequency.
\added{Satin weaves (e.g.\ 4- or 8-harness) have even fewer
interlacements, giving the best drape and lowest crimp, while
non-woven forms such as stitched non-crimp fabrics and randomly
oriented chopped-strand mat are chosen for straight, uncrimped fibres
and for low cost, respectively \citep{mallick2007,campbell2004}.} The selection
of architecture therefore involves balancing mechanical
requirements with manufacturing considerations.

Polymer matrices are broadly classified as thermosetting and
thermoplastic. Common thermosets include epoxy, polyester\added{, vinyl
ester} and phenolic resins, which form cross-linked networks during curing.
Thermoplastics, in contrast, soften upon heating and can be
reshaped\added{; common examples are polypropylene, polyamides and,
for high-temperature aerospace use, polyether ether ketone (PEEK)}. Epoxy-based prepregs contain fibres pre-impregnated
with a controlled resin quantity, improving material consistency
and eliminating separate resin application during lay-up.
Successful processing requires appropriate consolidation and
thermal curing to achieve the intended laminate properties
\citep{mallick2007,wu2026}.


\subsection{Laboratory Objectives}

This laboratory investigated the manufacture of a fibre-reinforced
polymer composite aerofoil using carbon- and glass-fibre prepregs
with a polyurethane foam core. The experiment involved manual
prepreg lay-up, vacuum bagging and thermal curing using an
aluminium mould \citep{chang2026}. The primary objectives were to:

\begin{itemize}[leftmargin=1.5em, itemsep=2pt, topsep=3pt]
    \item Understand the key stages of prepreg composite
    manufacturing, including material preparation, fibre
    lay-up, vacuum consolidation and curing.

    \item Examine the effects of fibre type, orientation,
    stacking sequence and processing conditions on
    composite quality and structural performance.

    \item Characterise the manufactured aerofoil through
    dimensional measurements, laminate architecture,
    fibre volume fraction estimation and surface inspection.

    \item Evaluate the advantages, limitations and
    engineering applications of manual composite lay-up
    relative to alternative manufacturing techniques.
\end{itemize}




    \section{Brief Theoretical Overview}
    % [Draft 1] The placeholder question list below is answered in
    % Section 3.1 and is kept here, commented out, for reference.
    % \begin{itemize}
    %     \item In the lab a simple aluminium billet mould was used. Outline the practical process of making a simple
    % mould of this type. In your report, please discuss the following points (this may need some research as
    % it contains content outside of the laboratory and is included to develop students’ self-guided research
    % skills):
    %     \begin{itemize}
    %         \item Female vs Male moulds
    %         \item  Differences in mould requirements for different composite-manufacturing techniques (i.e., wet-hand layup, vacuum infusion and prepreg (highlighting mould material selection)
    %         \item Common manufacturing methods for composite moulds.
    %         \item Use of gelcoats and other mould-surfacing chemicals and their purposes
    %     \end{itemize}
    % \end{itemize}

\begin{addedblock}
\subsection{Mould Design and Manufacture}
\label{sec:moulds}

\paragraph{Making a simple aluminium billet mould.}
A billet mould like the one used in the laboratory is machined from a
solid block of aluminium alloy \citep{campbell2004}. (i)~\emph{Design:}
the tool surface is modelled in CAD from the required outer surface of
the part, with a flat margin for trimming, sealant tape and bagging,
draft on near-vertical walls, and, for precision parts, an allowance
for thermal expansion at the cure temperature. (ii)~\emph{Stock
preparation:} a plate or billet (e.g.\ 6061-T6 or stress-relieved cast
tooling plate) is cut oversize and faced flat. (iii)~\emph{CNC
machining:} the cavity is roughed out with flat end mills and finished
with ball-nose cutters at a small step-over. (iv)~\emph{Finishing:}
the cusps are removed by wet sanding with progressively finer
abrasives and the surface is polished, since every tool defect is
reproduced on the part. (v)~\emph{Inspection and preparation:} the
profile is checked against the CAD model, and the surface is degreased,
sealed and given several cured coats of release agent.

\paragraph{Female versus male moulds.}
In a female (concave) mould the plies are laid inside a cavity, so the
mould-side face becomes the outer surface of the part and the external
dimensions are tightly controlled (Figure~\ref{fig:female_male}a).
This suits parts whose exterior is aerodynamically or cosmetically
critical, such as aerofoil skins, boat hulls and body panels, but
fabric can bridge across tight internal radii, leaving resin-rich or
poorly consolidated corners. In a male (convex) mould the plies are
laid over the tool, which controls the inner surface and dimensions
(Figure~\ref{fig:female_male}b), as needed for ducts or fairings that
fit over other structure. Lay-up is easier, but the outer face is the
rough bag side, which must be filled and finished if appearance
matters, and cure shrinkage makes the part grip the tool, so generous
draft is needed \citep{campbell2004,astrom1997}.
\end{addedblock}

\begin{figure}[htbp]
    \addedfloat
    \centering
    \begin{tikzpicture}[font=\footnotesize, >=Stealth, black]
        % ---- (a) female mould ----
        \fill[gray!30] (0,0) -- (6,0) -- (6,1.7) -- (5,1.7)
            arc[start angle=0, end angle=-180, x radius=2, y radius=1.1]
            -- (0,1.7) -- cycle;
        \draw[gray!60!black] (0,0) -- (6,0) -- (6,1.7) -- (5,1.7)
            arc[start angle=0, end angle=-180, x radius=2, y radius=1.1]
            -- (0,1.7) -- cycle;
        \draw[line width=2.4pt, black!85] (1.08,1.7)
            arc[start angle=180, end angle=360, x radius=1.92, y radius=1.02];
        \node at (3,0.22) {aluminium mould};
        \draw[->] (0.9,2.1) node[above, align=center] {bag side\\(rough)} -- (2.05,0.93);
        \draw[->] (5.0,2.1) node[above, align=center] {mould side\\(smooth outer\\surface)} -- (4.02,0.76);
        \node[below] at (3,-0.1) {(a) Female (concave) mould};
        % ---- (b) male mould ----
        \begin{scope}[xshift=7.6cm]
            \fill[gray!30] (0,0) -- (6,0) -- (6,0.6) -- (5,0.6)
                arc[start angle=0, end angle=180, x radius=2, y radius=1.1]
                -- (0,0.6) -- cycle;
            \draw[gray!60!black] (0,0) -- (6,0) -- (6,0.6) -- (5,0.6)
                arc[start angle=0, end angle=180, x radius=2, y radius=1.1]
                -- (0,0.6) -- cycle;
            \draw[line width=2.4pt, black!85] (5.08,0.6)
                arc[start angle=0, end angle=180, x radius=2.08, y radius=1.18];
            \node at (3,0.25) {aluminium mould};
            \draw[->] (0.9,2.1) node[above, align=center] {mould side\\(smooth inner\\surface)} -- (1.47,1.32);
            \draw[->] (5.1,2.1) node[above, align=center] {bag side\\(rough)} -- (4.45,1.6);
            \node[below] at (3,-0.1) {(b) Male (convex) mould};
        \end{scope}
    \end{tikzpicture}
    \caption{Schematic of a laminate (thick line) in (a) a female and
    (b) a male mould. The mould-side face of the part reproduces the
    tool finish.}
    \label{fig:female_male}
\end{figure}

\begin{addedblock}
\paragraph{Mould requirements for different processes.}
Table~\ref{tab:mould_req} compares the mould requirements of three
processes. Because wet hand lay-up cures near room temperature under
contact pressure only, inexpensive moulds of glass-fibre reinforced
plastic (GFRP), timber, MDF or plaster suffice. Vacuum infusion also
needs a vacuum-tight, stiff tool with a wide flange for sealant tape
and the resin and vacuum lines \citep{williams1996}. Prepreg moulds
must keep their shape and vacuum integrity through cures at typically
80--180\,$^\circ$C and, in an autoclave, up to about 7\,bar
\citep{campbell2004,hexcel2013}. Material selection for prepreg
tooling is driven by thermal expansion, thermal mass and cost.
Aluminium is light, cheap to machine and heats quickly, but its
coefficient of thermal expansion ($\alpha \approx
23\times10^{-6}$\,K$^{-1}$) far exceeds that of carbon/epoxy ($\alpha
\approx 2$--$3\times10^{-6}$\,K$^{-1}$). Heated by 100\,K to a typical
120\,$^\circ$C cure, an aluminium tool grows by about 0.23\,mm per
100\,mm, almost ten times more than the laminate, causing dimensional
error and tool--part interaction stresses that can distort thin parts
\citep{twigg2004}. Steel ($\alpha \approx 12\times10^{-6}$\,K$^{-1}$)
is durable but heavy, while Invar ($\alpha \approx
1.5\times10^{-6}$\,K$^{-1}$) and carbon/epoxy tooling closely match
carbon parts and are preferred for large, tight-tolerance aerospace
parts despite their cost \citep{campbell2004}. For a small aerofoil
like this one, aluminium offers the best balance of cost,
machinability and heat-up rate.
\end{addedblock}

\begin{table}[htbp]
    \addedfloat
    \centering
    \caption{Typical mould requirements for three composite
    manufacturing techniques \citep{campbell2004,astrom1997,williams1996,hexcel2013}.}
    \label{tab:mould_req}
    \footnotesize
    \renewcommand{\arraystretch}{1.15}
    \begin{tabularx}{\linewidth}{@{}>{\raggedright\arraybackslash}p{2.2cm}LLL@{}}
        \toprule
        & \textbf{Wet hand lay-up} & \textbf{Vacuum infusion}
        & \textbf{Prepreg (vacuum bag/oven or autoclave)} \\
        \midrule
        Cure temperature & Ambient (optional low-temperature post-cure)
            & Ambient to moderate & Typically 80--180\,$^\circ$C \\
        Pressure & Contact only (brush and roller)
            & Vacuum, at most $\approx$1\,atm
            & Vacuum; up to $\approx$7\,bar in an autoclave \\
        Key requirements & Smooth, sealed, releasable surface; low cost
            & Vacuum-tight and non-porous; stiff under vacuum; wide flange
            for sealant tape, resin inlet and vacuum line
            & Stable shape and vacuum integrity at cure temperature; CTE
            close to the part; low thermal mass; survives many thermal
            cycles \\
        Typical materials & GFRP with tooling gelcoat; timber, MDF,
            plaster or foam for one-offs & GFRP or epoxy composite,
            aluminium, glass or sheet steel & Aluminium, steel, Invar;
            carbon/epoxy or BMI tooling; high-temperature tooling board \\
        Surface preparation & Wax and/or PVA, or semi-permanent release;
            part gelcoat sprayed in mould & Semi-permanent release;
            gelcoat optional & Temperature-stable semi-permanent release;
            no gelcoat \\
        \bottomrule
    \end{tabularx}
\end{table}

\begin{addedblock}
\paragraph{Common mould manufacturing methods.}
(1)~Direct CNC machining of metal billet or plate (aluminium, steel,
Invar), as used here, gives accurate, durable, high-temperature tools
but becomes costly for large parts. (2)~Composite (splash) moulds are
taken from a master pattern, or plug, that is CNC-machined from
tooling board, MDF or foam (or 3D printed), then sealed and polished;
tooling gelcoat and a GFRP or CFRP tooling laminate are laid over the
plug and stiffened with a backing frame. This is the usual route for
large tools such as boat hulls and wind-turbine blades. (3)~Direct
machining of high-density polyurethane or epoxy tooling board suits
prototypes and short runs. (4)~Other routes include cast resin or cast
metal tools, electroformed nickel shells for very high surface quality
and additively manufactured moulds for rapid prototyping
\citep{campbell2004,astrom1997}.

\paragraph{Gelcoats and mould-surfacing chemicals.}
A gelcoat is a thixotropic, usually pigmented polyester, vinyl ester
or epoxy resin applied to the mould as the first layer (typically
about 0.5\,mm thick). It becomes the outer surface of the part,
providing colour and gloss, hiding fibre print-through and protecting
the laminate from UV, moisture and abrasion; tooling gelcoats are
harder, more heat-resistant grades that form the working surface of
composite moulds \citep{astrom1997}. Mould sealers fill the
microporosity of new or porous tool surfaces so that the release agent
forms a continuous film and resin cannot key into the tool. Release
agents then prevent adhesion: traditional carnauba waxes (several
buffed coats) and water-soluble polyvinyl alcohol (PVA) films, or
semi-permanent release agents that bond to the mould, tolerate high
cure temperatures and give several releases per application. Mould
cleaners remove resin residue and release-agent build-up between uses
\citep{explorecomposites_release}. In this laboratory only a liquid
release agent was wiped onto the bare aluminium moulds
(Section~\ref{sec:procedure}), as gelcoats are generally not used with
prepregs.

\subsection{Consumables}
\label{sec:consumables}
\end{addedblock}


The required consumables for composite manufacturing \added{vary} between processes. The consumables and their functions for the pre-preg layup performed in the laboratory are as follows:\\

\textbf{\added{Mould} Release Agent} is a coating applied to the \added{mould} to prevent the part from bonding to the \added{mould}. This ensures that \added{both} the part and \added{mould} are not damaged during the \added{demoulding} process.\\

A \textbf{Breather} is a layer of fabric that is placed on top of the layup and within the vacuum sealed bag. Its primary function is to provide a path for air to travel. Without it, the vacuum bag would suck itself flat to the part, resulting in pockets of air without direct exit paths. The breather creates a distributed pathway from all parts of the bag to ensure that all air is able to escape, producing a proper and effective vacuum.\\

\textbf{Peel Ply} is another layer of fabric that is laid atop and used to protect the outermost layer of a layup. It prevents the laminate from contamination by dust, oils, and other substances during handling, ensuring that the surface finish of the part is optimal. It also separates the part from the breather and bagging material, so that they can be easily removed. \added{When it is stripped after curing it leaves a clean, uniformly textured surface that is ideal for adhesive bonding; this mattered here because the bag-side surfaces of the skins were later bonded to the foam core \citep{campbell2004}.}\\

A \textbf{Vacuum Bag} is a flexible, airtight film bag. It surrounds and encloses the entire layup so that once air has been fully removed, atmospheric pressure \added{(up to about 101\,kPa)} compresses the layers, ensuring there are no voids/pockets and that the part is held together during the curing phase.\\

\textbf{Tacky Tape} is a sealant tape that is lined across the perimeter of the vacuum bag. The bag is pressed onto the tape to create the airtight seal necessary to create a vacuum.

\begin{addedblock}
\vspace{\baselineskip}
\textbf{Vacuum Port, Valve and Hose} connect the bag to the vacuum
pump, with a strip of breather keeping an open air path to the port.
More demanding prepreg lay-ups also use release film between the peel
ply and breather to control resin bleed, bleeder cloth to absorb excess
resin, heat-resistant flash tape, and thermocouples to monitor the part
temperature \citep{hexcel2013,campbell2004}. Figure~\ref{fig:bagging}
shows the complete arrangement.
\end{addedblock}

\begin{figure}[htbp]
    \addedfloat
    \centering
    \begin{tikzpicture}[font=\footnotesize, >=Stealth, black]
        % atmospheric pressure
        \foreach \x in {2.6,3.6,...,7.6} \draw[->, thick, gray!70!black] (\x,2.5) -- (\x,2.0);
        \node[above] at (5.1,2.5) {atmospheric pressure (up to $\approx$101\,kPa)};
        % bottom bag film
        \draw[blue!70!black, very thick] (0,0) -- (11.4,0);
        % mould
        \fill[gray!30] (1.4,0.05) rectangle (8.4,0.85);
        \draw[gray!60!black] (1.4,0.05) rectangle (8.4,0.85);
        \node at (5.1,0.42) {aluminium mould};
        % release agent
        \draw[orange!90!black, thick, dashed] (1.9,0.88) -- (7.9,0.88);
        % plies C/G/G/C
        \filldraw[fill=black!80, draw=black] (2.2,0.92) rectangle (7.6,1.04);
        \filldraw[fill=yellow!25, draw=black!60] (2.2,1.04) rectangle (7.6,1.16);
        \filldraw[fill=yellow!25, draw=black!60] (2.2,1.16) rectangle (7.6,1.28);
        \filldraw[fill=black!80, draw=black] (2.2,1.28) rectangle (7.6,1.40);
        % peel ply
        \filldraw[fill=orange!45, draw=orange!70!black] (2.1,1.43) rectangle (7.7,1.50);
        % breather, draped over the mould and led to the port
        \draw[line width=5pt, gray!45] (1.2,0.09) -- (1.2,1.63) -- (8.6,1.63) -- (8.6,0.09) -- (10.0,0.09);
        % top bag film
        \draw[blue!70!black, very thick, rounded corners=3pt] (0.35,0.14) -- (0.95,1.85) -- (8.85,1.85) -- (9.0,0.32) -- (10.95,0.32) -- (11.05,0.14);
        % tacky tape
        \filldraw[fill=yellow!70!orange, draw=black!60] (0.2,0.0) rectangle (0.5,0.14);
        \filldraw[fill=yellow!70!orange, draw=black!60] (10.9,0.0) rectangle (11.2,0.14);
        % vacuum port and hose
        \filldraw[fill=gray!70, draw=black] (9.25,0.19) rectangle (9.95,0.29);
        \filldraw[fill=gray!70, draw=black] (9.45,0.32) rectangle (9.75,0.75);
        \draw[line width=2pt, gray!60!black] (9.6,0.75) .. controls (9.6,1.2) and (10.0,1.25) .. (10.6,1.25);
        \node[right] at (10.6,1.25) {to pump};
        % labels
        \draw (2.15,1.47) -- (-0.4,1.95) node[left] {peel ply};
        \draw (2.25,1.16) -- (-0.4,1.45) node[left] {prepreg plies (C/G/G/C)};
        \draw (1.95,0.88) -- (-0.4,0.95) node[left] {release agent};
        \draw (0.3,0.07) -- (-0.4,0.4) node[left] {tacky tape};
        \draw (8.3,1.85) -- (9.3,2.3) node[right] {vacuum bag};
        \draw (8.6,1.2) -- (9.3,1.9) node[right] {breather};
        \draw (9.6,0.19) -- (9.6,-0.25) node[below] {vacuum port};
    \end{tikzpicture}
    \caption{Cross-section (not to scale) of the vacuum-bag arrangement
    used for each mould: the four prepreg plies on the release-coated
    aluminium mould are covered by peel ply and breather inside an
    envelope bag sealed with tacky tape and evacuated through the
    vacuum port.}
    \label{fig:bagging}
\end{figure}





\section{Lab Procedure}
\label{sec:procedure}

\textbf{Setup} - The following equipment was utilised to carry out the lab procedure: Disposable gloves, ruler, utility knife, female and male mould, pre-preg carbon fibre, pre-preg fibre glass, the consumables covered in section 3, and the vacuum port and line.\\

The procedure carried out in the lab is as follows:\\

\textbf{1} - Measure the \added{mould} dimensions using a ruler. Based on this measurement, cut the pre-preg carbon-fibre and fibre-glass to a suitable length and width, ensuring the cut covers the entire \added{mould}. There should be a total of four cuts \tbc{per mould?}, two carbon-fibre, and two fibre-glass.\\

\textbf{2} - Apply the releasing agent to the \added{mould}. Only a thin layer is necessary, however, ensure that it is evenly distributed without bare spots. This can be achieved by wiping it across the \added{mould} using a cloth.\\

\textbf{3} - Lay the carbon-fibre and fibre-glass onto the mould layer by layer. Ensure that each layer is pressed flat against the \added{mould} or layer beneath it, minimising the volume of air between layers. The order of the layers should create a symmetrical order about the mid-plane to ensure uniform properties once the part is cured \added{and to stop the laminate warping as it cools from the cure temperature}. This is achieved by laying \added{the} two carbon-fibre cuts on the top and bottom, sandwiching the two fibre-glass layers. The final result is a 4-ply composite.\\

\textbf{4} - Cut a layer of peel ply to the same dimensions as the fibre layups. Place this on top of the topmost layer, once again pressing it flat and minimising any trapped air.\\

\textbf{5} - With the main layup finished, place the \added{moulds} side by side on top of a vacuum bag. Apply tacky tape to the vacuum bag, creating a perimeter around the \added{moulds}.\\

\textbf{\added{6}} - Cut a single layer of breather fabric large enough to cover both \added{moulds}. Place the breather on top of the layup.\\

\textbf{\added{7}} - Seal the vacuum bag by pulling it back over the \added{moulds}. Gradually peel the backing paper of the tacky tape ensuring that the vacuum bag does not wrinkle when pressing it onto the tape, preventing small channels from emerging. Before finalising the seal, leave one edge of the bag open.\\

\textbf{\added{8}} - Place the vacuum port at a distance between the \added{mould} and \added{bag's} edge. The base should sit inside the bag, and screw the top outside the bag, cutting a small slit to allow air to escape. To ensure air has an undisturbed pathway to exit via the port, cut a strip of breather fabric that leads from the primary breather layer to the vacuum port.\\

\textbf{\added{9}} - Seal the final edge of the vacuum bag. To test the quality of the vacuum, attach the hose to the port and begin the air evacuation process. If there is a noticeable hissing sound, it is highly likely that air is escaping from \added{an} unsealed hole around the vacuum bag perimeter.\\

\textbf{\added{10}} - Once the seal quality is verified, the entire assembly is transferred into the curing oven. The vacuum hose line is connected to a port within the oven. The program should be set for optimum curing cycle \added{recommended by the prepreg manufacturer} \tbc{state the cure cycle used: heating rate, cure temperature and hold time}. The oven is then started and fan turned on for uniform heating.\\

\textbf{\added{11}} - Once the curing cycle has been completed and the part has cooled sufficiently, it can be removed from the oven. Once the vacuum is vented, the bag can be cut open and the \added{moulds} removed from within. Consumables (\added{e.g.\ breather}, peel ply) can be removed. The part should snap out of the \added{mould} with little problem due to the use of the release agent.\\

\textbf{\added{12}} -  The result should be two fibre-reinforced composite parts that make up the upper and lower surface of an \added{aerofoil}. These parts are then placed on the outer surface of a foam core with the intended \added{aerofoil} profile and bonded using a suitable adhesive \tbc{name the adhesive used}.



\section{Results}

\begin{addedblock}
\subsection{Manufactured Aerofoil}

The finished part (Figure~\ref{fig:part}) is an aerofoil section made
of two cured hybrid skins, forming the upper and lower surfaces,
bonded to a polyurethane foam core. It has a chord of \tbc{}\,mm, a
span of \tbc{}\,mm, a maximum thickness of \tbc{}\,mm ($t/c =$
\tbc{}\,\%) and a total mass of \tbc{}\,g. Skin thickness was measured
with \tbc{a micrometer, $\pm$0.01\,mm} at \tbc{n} locations on each
skin away from the trimmed edges, lengths with \tbc{vernier calipers,
$\pm$0.02\,mm} and masses with \tbc{a digital balance, $\pm$0.01\,g}.
\end{addedblock}

\begin{figure}[htbp]
    \addedfloat
    \centering
    \begin{minipage}[c]{0.34\linewidth}
        \centering
        \fbox{\parbox[c][3.4cm][c]{0.92\linewidth}{\centering
            \tbc{insert photograph of the finished aerofoil (mould-side
            and bag-side views)}}}\\[2pt]
        {\footnotesize (a)}
    \end{minipage}\hfill
    \begin{minipage}[c]{0.64\linewidth}
        \centering
        \begin{tikzpicture}[font=\footnotesize, >=Stealth, black,
            declare function={xc(\t)=(1-cos(\t))/2;
                yt(\x)=5*0.15*(0.2969*sqrt(\x)-0.1260*\x-0.3516*\x^2+0.2843*\x^3-0.1015*\x^4);}]
            \def\c{5.6}
            % foam core and skins (symmetric 15 % section, schematic only)
            \fill[yellow!25] plot[domain=0:180, samples=70, variable=\t]
                ({\c*xc(\t)}, {\c*yt(xc(\t))}) --
                plot[domain=180:0, samples=70, variable=\t]
                ({\c*xc(\t)}, {-\c*yt(xc(\t))}) -- cycle;
            \draw[line width=1.6pt, black!85] plot[domain=0:180, samples=70, variable=\t]
                ({\c*xc(\t)}, {\c*yt(xc(\t))});
            \draw[line width=1.6pt, black!60] plot[domain=0:180, samples=70, variable=\t]
                ({\c*xc(\t)}, {-\c*yt(xc(\t))});
            \node at (2.2,0) {PU foam core};
            \draw[->] (1.0,1.0) node[above] {upper skin} -- (1.0,0.47);
            \draw[->] (1.0,-1.0) node[below] {lower skin} -- (1.0,-0.47);
            \draw[<->] (0,-1.45) -- node[below] {chord $c$} (\c,-1.45);
            % ply stack detail
            \draw[dashed, gray] (3.0,0.43) -- (5.9,0.75);
            \begin{scope}[shift={(5.9,0.0)}]
                \filldraw[fill=black!80, draw=black] (0,0.9) rectangle (0.9,1.2);
                \filldraw[fill=yellow!25, draw=black!60] (0,0.6) rectangle (0.9,0.9);
                \filldraw[fill=yellow!25, draw=black!60] (0,0.3) rectangle (0.9,0.6);
                \filldraw[fill=black!80, draw=black] (0,0.0) rectangle (0.9,0.3);
                \node[right, font=\scriptsize] at (0.9,1.05) {4 carbon (bag side)};
                \node[right, font=\scriptsize] at (0.9,0.75) {3 E-glass};
                \node[right, font=\scriptsize] at (0.9,0.45) {2 E-glass};
                \node[right, font=\scriptsize] at (0.9,0.15) {1 carbon (mould side)};
                \draw[decorate, decoration={brace, amplitude=3pt}] (-0.08,0) -- node[left=3pt] {$h$} (-0.08,1.2);
            \end{scope}
        \end{tikzpicture}\\[2pt]
        {\footnotesize (b)}
    \end{minipage}
    \caption{(a) The manufactured aerofoil. (b) Schematic section (not
    to scale) showing the two skins bonded to the foam core and the
    $[\mathrm{C}/\mathrm{G}]_\mathrm{s}$ ply sequence of each skin.}
    \label{fig:part}
\end{figure}

\begin{addedblock}
\subsection{Laminate Architecture}

Each skin is a four-ply, symmetric, inter-ply hybrid laminate with a
woven carbon/epoxy ply on each face and two woven E-glass/epoxy plies
at the mid-plane (Table~\ref{tab:architecture}). With the warp of
every fabric aligned with the span, the lay-up is
$[0_\mathrm{C}/0_\mathrm{G}]_\mathrm{s}$, where each 0 denotes a
0/90$^\circ$ woven ply \tbc{confirm ply angles, fabric styles and
prepreg products}. Carbon and glass each account for half of the plies
and, because their nominal areal masses are equal, half of the fibre
mass. However, carbon fibre (T300-type, 1.76\,g\,cm$^{-3}$
\citep{toray_t300}) is about 30\,\% less dense than E-glass
(2.54\,g\,cm$^{-3}$ \citep{mallick2007}), so carbon makes up about
59\,\% of the fibre volume. The cured epoxy matrix has a density of
about 1.2\,g\,cm$^{-3}$ \citep{mallick2007}.
\end{addedblock}

\begin{table}[htbp]
    \addedfloat
    \centering
    \caption{Laminate architecture of each skin ($A_w$ = fabric areal
    mass, $\rho_f$ = fibre density, $W_r$ = prepreg resin content by
    mass; matrix density $\rho_m = 1.2$\,g\,cm$^{-3}$). Values marked
    $^*$ are nominal values typical of woven epoxy prepregs, used until
    the supplier data sheets are checked \tbc{replace with data-sheet
    values}.}
    \label{tab:architecture}
    \small
    \renewcommand{\arraystretch}{1.1}
    \begin{tabular}{@{}cllcccc@{}}
        \toprule
        \textbf{Ply} & \textbf{Position} & \textbf{Prepreg / fabric}
        & \textbf{Angle} & $A_w$ (g\,m$^{-2}$) & $\rho_f$ (g\,cm$^{-3}$)
        & $W_r$ (wt\%) \\
        \midrule
        1 & Mould side & Carbon/epoxy, woven \tbc{weave} & 0/90 & 200$^*$ & 1.76 & 40$^*$ \\
        2 & Core & E-glass/epoxy, woven \tbc{weave} & 0/90 & 200$^*$ & 2.54 & 35$^*$ \\
        3 & Core & E-glass/epoxy, woven & 0/90 & 200$^*$ & 2.54 & 35$^*$ \\
        4 & Bag side & Carbon/epoxy, woven & 0/90 & 200$^*$ & 1.76 & 40$^*$ \\
        \bottomrule
    \end{tabular}
\end{table}

\begin{addedblock}
\subsection{Fibre Volume Fraction}

Because the fabric areal masses are known, the fibre volume fraction
can be estimated non-destructively from the cured thickness, using
the areal-weight approach of ASTM D3171 Method~II \citep{astm_d3171}:
\begin{equation}
    V_f = \frac{1}{h}\sum_i \frac{n_i A_{w,i}}{\rho_{f,i}},
    \label{eq:vf}
\end{equation}
where $n_i$ is the number of plies of fabric $i$ and $h$ the measured
laminate thickness. For a void-free laminate with no resin bleed, the
cured thickness of each ply follows from the prepreg resin content:
\begin{equation}
    t_{\mathrm{ply},i} = \frac{A_{w,i}}{\rho_{f,i}}
    + \frac{A_{w,i}}{\rho_m}\,\frac{W_{r,i}}{1-W_{r,i}}.
    \label{eq:cpt}
\end{equation}
With the values in Table~\ref{tab:architecture}, Eq.~(\ref{eq:cpt})
gives cured ply thicknesses of 0.225\,mm for carbon ($V_f=50.6$\,\%)
and 0.168\,mm for glass ($V_f=46.7$\,\%), so the expected skin
thickness is $h_\mathrm{th}=0.786$\,mm with $V_f\approx48.9$\,\%
(28.9\,\% carbon, 20.0\,\% glass and 51.1\,\% matrix by volume), an
areal mass of 1.28\,kg\,m$^{-2}$ and a density of 1.63\,g\,cm$^{-3}$.
The summed fibre thickness $\sum n_i A_{w,i}/\rho_{f,i}$ is 0.385\,mm,
so Eq.~(\ref{eq:vf}) reduces to $V_f = 0.385/h$ ($h$ in mm), and
Table~\ref{tab:vf} compares the measured skins with this nominal case.
Since $V_f \propto 1/h$, a thickness error of only 0.05\,mm shifts
$V_f$ by about 3 percentage points, so the mean of several readings
should be used. If the skin masses and areas are known, comparing the
measured density with the void-free value (1.63\,g\,cm$^{-3}$ at the
nominal resin content) also gives an estimate of the void content
\citep{astm_d3171}.
\end{addedblock}

\begin{table}[htbp]
    \addedfloat
    \centering
    \caption{Measured skin thickness and fibre volume fraction from
    Eq.~(\ref{eq:vf}), compared with the nominal void-free laminate.}
    \label{tab:vf}
    \small
    \renewcommand{\arraystretch}{1.1}
    \begin{tabular}{@{}lccc@{}}
        \toprule
        \textbf{Quantity} & \textbf{Nominal (void-free)} & \textbf{Upper skin} & \textbf{Lower skin} \\
        \midrule
        Thickness $h$, mean $\pm$ s.d.\ (mm) & 0.786 & \tbc{} & \tbc{} \\
        Mass (g) & -- & \tbc{} & \tbc{} \\
        Fibre volume fraction $V_f$ (\%) & 48.9 & \tbc{} & \tbc{} \\
        Carbon / glass fibre volume (\%) & 28.9 / 20.0 & \tbc{} & \tbc{} \\
        Matrix + voids (\%) & 51.1 & \tbc{} & \tbc{} \\
        \bottomrule
    \end{tabular}
\end{table}

\begin{addedblock}
\subsection{Surface Quality}

\tbc{replace the bracketed observations below with what was seen on
your part.} The mould-side (aerodynamic) surfaces were \tbc{e.g.\
smooth and glossy, reproducing the aluminium finish, with the carbon
weave clearly visible}. The bag-side surfaces carried the uniform matte
texture imprinted by the peel ply, which gives a clean, keyed surface
for bonding to the core. The defects observed were \tbc{e.g.\
pinholes or porosity, dry or resin-starved patches, resin-rich areas,
wrinkles near the leading or trailing edge, fraying at trimmed edges,
gaps or steps at the skin--core joints}, located \tbc{where}. After
demoulding the skins \tbc{did / did not} show warping.
\end{addedblock}

\section{Discussion}

\begin{addedblock}
\subsection{Quality of the Manufactured Part}

\tbc{compare the measured thickness and $V_f$ with the nominal values
in Table~\ref{tab:vf}.} A skin thicker than the nominal 0.786\,mm
($V_f$ below 49\,\%) points to trapped air, resin-rich regions or
incomplete consolidation, whereas a thinner skin indicates resin bleed
into the peel ply and breather. Several features of this process
favour some porosity: consolidation relied on vacuum alone (at most
$\approx$1\,atm), the plies were not debulked individually
\tbc{confirm}, and air trapped between plies of conventional prepreg
can escape only through the laminate edges, whereas an autoclave adds
several bar of pressure that collapses voids
\citep{centea2015,hexcel2013}. Manual cutting and placement also cause
ply misalignment of a few degrees, wrinkles where the fabric is forced
around the tight leading-edge curvature, and bridging across concave
regions of the moulds; together with the adhesive bond line, these
explain differences between the measured and nominal profile. Quality
could be improved by debulking under vacuum after each ply, cutting
plies from templates with marked fibre directions, leak-testing the
bag before cure, adding a low-temperature dwell to the cure cycle so
that air can escape before the resin gels, and, for critical parts,
curing in an autoclave \citep{hexcel2013,centea2015}.

\subsection{Comparison of Manufacturing Methods}

FRP manufacturing processes trade capital cost, labour and production
rate against laminate quality (Table~\ref{tab:processes}). Open-mould
processes (wet hand lay-up and spray-up) need the least investment but
give low fibre fractions and high, operator-dependent void contents.
Vacuum infusion and resin transfer moulding (RTM) make low-void,
higher-$V_f$ parts from dry fabric: infusion needs only a single-sided
tool and suits very large parts \citep{williams1996}, while the matched
tools of RTM give two finished surfaces and shorter cycles for medium
volumes. Prepreg processes give the most precise control of fibre
content and orientation. Autoclave curing, including automated tape
laying and fibre placement, gives the highest quality for aerospace
primary structure, and the vacuum-bag/oven route used here achieves
much of that quality at far lower capital cost, especially with
prepregs formulated for out-of-autoclave cure \citep{centea2015}.
Filament winding and pultrusion reach high $V_f$ and production rates
but are limited to bodies of revolution and constant cross-sections,
respectively \citep{campbell2004,astrom1997,mallick2007}.
\end{addedblock}

\begin{table}[htbp]
    \addedfloat
    \centering
    \caption{Comparison of common FRP manufacturing processes.
    $V_f$ values are indicative only and depend strongly on materials
    and practice \citep{campbell2004,astrom1997,mallick2007,williams1996}.}
    \label{tab:processes}
    \footnotesize
    \renewcommand{\arraystretch}{1.1}
    \begin{tabularx}{\linewidth}{@{}>{\raggedright\arraybackslash}p{2.6cm}Lcc>{\raggedright\arraybackslash}p{1.7cm}>{\raggedright\arraybackslash}p{3.6cm}@{}}
        \toprule
        \textbf{Process} & \textbf{Principle} & $\boldsymbol{V_f}$ \textbf{(\%)}
        & \textbf{Voids} & \textbf{Tooling / capital cost} & \textbf{Volume; typical applications} \\
        \midrule
        Wet hand lay-up & Resin worked into dry fabric by brush and roller
            in an open mould & 20--40 & High & Low / low
            & One-off to low; boat hulls, prototypes, repairs \\
        Spray-up & Chopped roving and resin sprayed onto an open mould
            & 15--30 & High & Low / low & Low--medium; baths, tanks,
            non-structural panels \\
        Prepreg, vacuum bag and oven (this lab) & Prepreg plies laid up,
            vacuum bagged and oven cured & 50--60 & Low--mod.
            & Medium / medium & Low--medium; UAVs, motorsport, sporting goods \\
        Prepreg, autoclave (incl.\ ATL/AFP) & As above, cured under heat
            and 3--7\,bar & 55--65 & Very low & High / very high
            & Low--medium; aerospace primary structure \\
        Vacuum infusion & Resin drawn through a dry preform under a bag
            by vacuum & 50--60 & Low & Low / low & Low--medium; wind
            blades, boat hulls \\
        Resin transfer moulding & Resin injected into a dry preform in a
            closed matched mould & 40--60 & Low & High / medium
            & Medium--high; automotive and aerospace parts \\
        Compression moulding (SMC) & Chopped-fibre charge pressed in
            heated metal dies & 15--35 & Low & Very high / high
            & High; automotive body panels \\
        Filament winding & Resin-wetted tows wound onto a rotating mandrel
            & 55--70 & Low & Medium / medium & Medium--high; pressure
            vessels, pipes, shafts \\
        Pultrusion & Fibres pulled through a resin bath and heated die
            & 50--70 & Low & Medium / medium & Continuous; constant-section
            profiles \\
        \bottomrule
    \end{tabularx}
\end{table}

\begin{addedblock}
\subsection{Parameters Influencing Composite Properties}

\paragraph{Constituents.} Fibre type sets the ceiling on properties:
T300-type carbon fibre ($E_f=230$\,GPa, $\rho_f=1.76$\,g\,cm$^{-3}$) is
about 4.6 times stiffer per unit mass than E-glass ($E_f\approx72$\,GPa,
$\rho_f=2.54$\,g\,cm$^{-3}$) \citep{toray_t300,mallick2007}, whereas
glass is cheaper, has a higher strain to failure and is an electrical
insulator. Hybridising them, as here, trades some stiffness for lower
cost and better damage tolerance \citep{swolfs2014}. The matrix
governs the matrix-dominated properties (transverse, shear,
compressive and interlaminar strength), the maximum service
temperature (through its glass-transition temperature) and
environmental resistance, while the fibre--matrix interface controls
load transfer \citep{mallick2007}.

\paragraph{Fibre volume fraction.} Fibre-direction stiffness follows
the rule of mixtures, $E_1 = V_f E_f + (1-V_f)E_m$
\citep{mallick2007}, so each extra 1\,\% of carbon fibre raises the
modulus of a unidirectional carbon/epoxy ply by about 2.3\,GPa. In
practice $V_f$ is limited to about 60--70\,\% by fibre packing and
wet-out, and it is set mainly by the consolidation pressure and resin
bleed of the process (Table~\ref{tab:processes}).

\paragraph{Fibre orientation and stacking sequence.} Properties are
highest along the fibres and fall rapidly off-axis: Krenchel's
orientation efficiency factor is 1 for unidirectional fibres, 0.5 for
a 0/90$^\circ$ fabric loaded along one fibre direction and 0.375 for
random in-plane mat \citep{krenchel1964}. Woven 0/90$^\circ$ plies
therefore give balanced properties in two directions but low in-plane
shear stiffness, so $\pm45^\circ$ plies are added where torsional
stiffness is needed, as in wing skins. A symmetric lay-up removes
bending--extension coupling, so the skins do not warp when the carbon
and glass plies contract by different amounts on cooling from the cure
temperature \citep{mallick2007}. Placing the stiffest plies furthest
from the mid-plane maximises bending stiffness: a simple lamination
estimate using the nominal ply thicknesses and effective ply moduli of
about 60\,GPa (carbon) and 19\,GPa (glass) shows that the
$[\mathrm{C}/\mathrm{G}]_\mathrm{s}$ skin has about 2.1 times the
flexural rigidity of a $[\mathrm{G}/\mathrm{C}]_\mathrm{s}$ skin with
the same materials and in-plane stiffness.

\paragraph{Void content.} Voids act as stress concentrators and crack
initiation sites and mainly degrade matrix-dominated properties:
interlaminar shear strength falls by roughly 7\,\% for each 1\,\% of
voids, up to about 4\,\% void content \citep{judd1978}. Voids come from
air trapped during lay-up, moisture and volatiles, and insufficient
consolidation pressure, which is why debulking, vacuum integrity and
pressure are critical \citep{centea2015}.

\paragraph{Processing conditions.} The cure cycle (heating rate, dwell
temperature and time) controls resin viscosity and flow, the degree of
cure and hence the glass-transition temperature, and the exotherm in
thick sections, while prepreg that has exceeded its out-life loses
tack and flow and becomes more porous \citep{hexcel2013}. The
manufacturing method sets the achievable $V_f$, void content and fibre
straightness (Table~\ref{tab:processes}), and manual operations add
ply misalignment and wrinkles, so nominally identical laminates can
differ noticeably.

\subsection{Application of Hand Lay-up}

Hand lay-up suits the wing skins of small fixed-wing unmanned aerial
vehicles (UAVs) built in batches of a few to a few hundred by
start-ups, research groups or student teams. Such volumes cannot
justify the matched tooling of RTM or an autoclave, whereas hand
lay-up needs only a single-sided mould (machined aluminium, tooling
board or even 3D printed), a vacuum pump and an oven. Female moulds
give the smooth, accurate outer surface needed for low drag, as this
laboratory showed. Because every ply is placed by hand, the laminate
can be tailored locally, with unidirectional carbon spar caps for
bending, $\pm45^\circ$ plies for torsion and extra plies at attachment
points, and glass can replace carbon around antennas because, unlike
carbon, it does not block radio signals. Design changes need only a
new mould or ply schedule, allowing rapid iteration, and the labour
cost and operator-dependent quality are acceptable at these volumes,
particularly when prepreg and vacuum bagging are used as here.
\end{addedblock}

\section{Conclusion}

\begin{addedblock}
An FRP aerofoil was made by hand lay-up of woven carbon and E-glass
epoxy prepregs in aluminium billet moulds, vacuum bagging and oven
curing, and the skins were bonded to a polyurethane foam core.
\begin{itemize}[leftmargin=1.5em, itemsep=1pt, topsep=3pt]
    \item Each skin is a symmetric four-ply hybrid
    $[\mathrm{C}/\mathrm{G}]_\mathrm{s}$ laminate in which carbon is half
    of the fibre mass but, being less dense (1.76 versus
    2.54\,g\,cm$^{-3}$), about 59\,\% of the fibre volume.
    \item The fibre volume fraction from the measured thickness was
    \tbc{}\,\%, against about 49\,\% for a void-free laminate,
    indicating \tbc{voids / resin bleed}. The mould-side surfaces
    \tbc{summary}, and \tbc{main defects}.
    \item Moulds must suit the process: aluminium is economical and
    heats quickly, but expands almost ten times more than carbon/epoxy,
    so Invar or composite tooling is preferred for large precision parts.
    \item Fibre type, volume fraction, orientation, stacking sequence,
    void content and cure govern laminate properties. Placing carbon
    outermost roughly doubles the skin's flexural rigidity, and the
    symmetric lay-up prevents warping.
    \item Hand lay-up remains attractive for low-volume, high-value parts
    such as UAV wing skins; ply-by-ply debulking or autoclave pressure
    would further reduce porosity.
\end{itemize}
\end{addedblock}

 \bibliographystyle{unsrtnat}
\bibliography{references}

    \section{Appendix}
\tbc{the Appendix is currently empty. Delete this heading, or add raw
measurement data here; note that appendix pages count towards the
10-page limit.}

    \end{document}
